\section{Conclusion}
\label{sec:conclusion}

We formulated continual PEFT as online allocation within a fixed low-rank
update space under task-identity-free, clocked temporal mixture drift. The resulting analysis
does not equate forgetting with capacity. It separates irreducible conflict,
fixed-rank approximation, and online error; characterizes when sequential
updates accumulate old loss; and shows that the locally zero-interference space
contracts as historical sensitivity directions accumulate. This geometry
motivates \method{}, a single fixed-rank adapter that, each window, re-solves a
risk-reweighted reduced-rank regression from a bounded buffer rather than growing
or freezing capacity. We are explicit about what allocation can and cannot do:
the spectral tail is an unavoidable capacity floor that no allocation rule
removes, so reallocation can only reduce the avoidable online error and steer
where the fixed budget spends its retention. In our controlled linear mixture
every implicit allocator charges that aggregate floor to the worst-served
source, so their worst-group retention clusters at a common low value, and only
the coverage-aware explicit minimax re-solve \emph{redistributes} the fixed
budget to raise it---at matched frequency-weighted mean, as a paired-average gain
over a crossed design of held-out geometries and stream draws, not a lowering of
the aggregate floor and not a uniform reduction in forgetting. A minimal
nonlinear probe (a two-layer ReLU MLP) does not yet reproduce the
linear-regime prediction, which we report rather than suppress. The next decisive step is
empirical: determine whether the predicted spectral and curvature quantities
forecast forgetting on real temporal streams and whether curvature-aware
allocation improves a clearly specified retention--plasticity or worst-group
trade-off under strictly matched resources.
